\documentclass[11pt,aspectratio=169]{beamer}
\usetheme{Boadilla}
\usefonttheme{professionalfonts}
\setbeamertemplate{navigation symbols}{}

\usepackage[T1]{fontenc}
\usepackage{lmodern}
\usepackage{amsmath,amssymb,mathtools}
\usepackage{booktabs}
\usepackage{tabularx}
\usepackage{array}
\usepackage{graphicx}
\usepackage[percent]{overpic}
\usepackage{xcolor}
\usepackage[authoryear,round]{natbib}
\usepackage{bibentry}
\bibliographystyle{apalike}
\setbeamertemplate{bibliography item}{}
\makeatletter
\def\hyper@natlinkstart#1{}
\def\hyper@natlinkend{}
\def\hyper@natlinkbreak#1#2{#1}
\makeatother

\graphicspath{{figs/}}

% Section headline for the audience arc.
\setbeamertemplate{headline}{%
  \hbox{}\vskip2pt\hfill\textcolor{gray!70}{\scriptsize\insertsectionhead}\hspace{1em}\vskip-6pt%
}
\setbeamertemplate{section in toc}[sections numbered]
\setbeamertemplate{frametitle}[default][left]
\setbeamerfont{framesubtitle}{size=\small}

\newcommand{\figplaceholder}[1]{\begin{center}\fbox{\begin{minipage}[c][0.32\textheight][c]{0.82\textwidth}\centering\large #1\end{minipage}}\end{center}}
\newcommand{\shortfigplaceholder}[1]{\begin{center}\fbox{\begin{minipage}[c][0.22\textheight][c]{0.82\textwidth}\centering\large #1\end{minipage}}\end{center}}
\newcommand{\compactfigplaceholder}[1]{\begin{center}\fbox{\begin{minipage}[c][0.10\textheight][c]{0.82\textwidth}\centering\normalsize #1\end{minipage}}\end{center}}
\newcommand{\R}{\mathcal{R}}
\newcommand{\E}{\mathcal{E}}
\newcommand{\dd}{\,\mathrm{d}}
\newcommand{\good}[1]{\textcolor{blue!70!black}{#1}}
\newcommand{\bad}[1]{\textcolor{red!70!black}{#1}}
\newcommand{\SAR}{speed asymmetry ratio}
% Inline slide navigation (end of sentence / bullet).
\newcommand{\app}[1]{\allowbreak\hspace{0.25em}\hyperlink{#1}{\beamergotobutton{details}}}
\newcommand{\back}[1]{\allowbreak\hspace{0.25em}\hyperlink{#1}{\beamerreturnbutton{back}}}
% Appendix / frame footer: bottom-right, not centered.
\newcommand{\navback}[1]{\par\vspace{0.35em}\noindent\hfill\hyperlink{#1}{\beamerreturnbutton{back}}}
\newcommand{\navfwd}[1]{\par\vspace{0.35em}\noindent\hfill\hyperlink{#1}{\beamergotobutton{details}}}
\newcommand{\navbackfwd}[2]{\par\vspace{0.35em}\noindent\hfill\hyperlink{#1}{\beamerreturnbutton{back}}\;\hyperlink{#2}{\beamergotobutton{details}}}

\title[Directional Congestion]{Directional Congestion}
\subtitle{Measurement, Spatial Equilibrium, and Tidal-Lane Reallocation}
\author[Duan, Feng, Gu]{Yige Duan (SJTU) $\quad$ Xiaoruo Feng (PKU) $\quad$ Yizhen Gu (PKU)}
\date[SMUE Wuhan, 2026]{SMUE Wuhan, 2026}

\begin{document}

\begin{frame}
\titlepage
\end{frame}

% =====================================================================
% Opener: roadmap followed by examples of directional congestion.
% =====================================================================
\begin{frame}{Roadmap}
\footnotesize\tableofcontents
\end{frame}

\begin{frame}{Directional Congestion in Cities}

\centering
% All four photos center-cropped to 3:2 for a uniform grid.
\begin{minipage}[t]{0.32\textwidth}\centering
\includegraphics[width=0.8\linewidth]{cityBeijing_c.jpg}\\[0.02em]{\scriptsize Beijing}
\end{minipage}\hspace{0.02\textwidth}%
\begin{minipage}[t]{0.32\textwidth}\centering
\includegraphics[width=0.8\linewidth]{cityWuhan_c.jpg}\\[0.02em]{\scriptsize Wuhan}
\end{minipage}

\vspace{0.25em}
\begin{minipage}[t]{0.32\textwidth}\centering
\includegraphics[width=0.8\linewidth]{cityShenzhen2_c.jpg}\\[0.02em]{\scriptsize Shenzhen}
\end{minipage}\hspace{0.02\textwidth}%
\begin{minipage}[t]{0.32\textwidth}\centering
\includegraphics[width=0.8\linewidth]{cityKualaLumpur_c.jpg}\\[0.02em]{\scriptsize Kuala Lumpur}
\end{minipage}
\end{frame}

% =====================================================================
\section{Introduction}
% =====================================================================
\begin{frame}{Motivation}
\begin{columns}[T,onlytextwidth]
\column{0.65\textwidth}
\begin{itemize}\small
  \item \textbf{Directional congestion is widely observed on urban roads, particularly during peak hours.}
  \begin{itemize}
    \item Commuting traffic is shaped by where workers live and where jobs are located.
  \end{itemize}
  \item \textbf{The literature often ignores this asymmetry in congestion.}
  \begin{itemize}
    \item Examples: QSM applications usually enter travel frictions through bilateral or route-level costs \citep{ahlfeldt2015,allenarkolakis2022}.
    \item Except: transport studies discuss peak-direction traffic \citep{vickrey1969,wolshon2006}.
  \end{itemize}
  \item \textbf{Tidal lanes are proposed as a low-cost way to address it.}
  \begin{itemize}
    \item However, little is known about the short-run and long-run effects of tidal lanes.
  \end{itemize}
\end{itemize}
\column{0.34\textwidth}
\centering
\href{run:./A6_XiaoruoFeng_DirectionalCongestion_video.mov}{%
  \includegraphics[width=\linewidth,height=0.70\textheight,keepaspectratio]{TidalVedio/pic.jpg}}\\[3pt]
{\scriptsize \good{Click to play}}
\end{columns}
\end{frame}

\begin{frame}{This Paper}
\small
\begin{itemize}
  \item \textbf{Measure directionality, using speed data at both the trip and the road segment level.}
  \begin{itemize}\footnotesize
    \item National trip-level: flow-weighted median \SAR\ is \good{0.878}.
    \item Within-city road segments: \SAR\ $\leq1/3$ flags \good{179} Beijing AM two-way centerlines (\good{4.9\%}, about \good{213 km}).
  \end{itemize}
  \item \textbf{Quantify what symmetry misses, using a spatial GE model.}
  \begin{itemize}\footnotesize
    \item Holding fundamentals and parameters fixed, replacing directed costs with reciprocal symmetric costs changes $\bar W$ by \good{$1.10\%$} and moves \good{$3.36\%$} of residents and \good{$6.15\%$} of jobs.
  \end{itemize}
  \item \textbf{Evaluate counterfactuals under endogenous congestion GE.}
  \begin{itemize}\footnotesize
    \item Tidal top-100 shifts one lane toward the slow AM direction on each treated corridor.
    \item Welfare rises \good{$+0.87\%$} in ST and \good{$+0.32\%$} in LT.
    \item Capacity-expansion benchmark builds an additional \good{$1{,}100$ lane-km}, while tidal lanes add \good{zero}.
  \end{itemize}
\end{itemize}
\end{frame}

\begin{frame}{Literature}
\begin{itemize}\small
  \item \textbf{Directional congestion}
  \begin{itemize}\footnotesize
    \item Peak-period bottlenecks, measured urban speeds, and reversible-lane operations \citep{vickrey1969,arnott1993,couture2018,akbar2023,wolshon2006}.
    \item \good{This paper measures directionality at national trip level and road-segment level using high-frequency speed data.}
  \end{itemize}
  \item \textbf{Spatial equilibrium}
  \begin{itemize}\footnotesize
    \item Economic geography and quantitative spatial models \citep{krugman1991,eaton2002,allen2014,ahlfeldt2015}.
    \item \good{This paper applies spatial GE with directional congestion and analyze tidal lane counterfactuals}
  \end{itemize}
  \item \textbf{Transportation policy and congestion}
  \begin{itemize}\footnotesize
    \item Demand-side policies: congestion pricing and equilibrium sorting \citep{vickrey1969,kreindler2024,barwick2024,yang2020}.
    \item Supply-side policies: road capacity, fundamental law, and optimal networks \citep{duranton2011,fajgelbaum2020}.
    \item Public transit as a substitute: subways relieve road congestion \citep{gu2021}.
    \item \good{This paper evaluates tidal lanes as a directed capacity-reallocation policy.}
  \end{itemize}
\end{itemize}
\end{frame}

\begin{frame}[label=nattrip]{Directionality Appears at National and Road-Segment Scales}
\begin{itemize}\footnotesize
  \item \textbf{Definition.} For each pair, the \SAR\ is the lower speed divided by the higher speed across its two directions. It lies between 0 and 1; lower values indicate stronger directional asymmetry.
\end{itemize}
\vspace{-0.35em}
\begin{columns}[T,onlytextwidth]
\column{0.66\textwidth}
\centering\scriptsize \textbf{National trip-level directionality}\par
\vspace{0.15em}
\begin{minipage}[t]{0.49\linewidth}\centering
\includegraphics[width=\linewidth,height=0.31\textheight,keepaspectratio]{pat_national_ratio_hist.png}
\end{minipage}\hspace{0.01\linewidth}%
\begin{minipage}[t]{0.49\linewidth}\centering
\includegraphics[width=\linewidth,height=0.31\textheight,keepaspectratio]{pat_congestion_am_scatter.png}
\end{minipage}
\vspace{-0.35em}
\begin{itemize}\scriptsize
  \item Flow-weighted median \SAR\ is \good{0.878} across the top 100 congested cities.
  \item \good{58\%} of pair-flow faces \SAR\ $\leq0.90$ during the AM peak.
  \item More congested cities tend to have lower ratios and stronger asymmetry ($r=\bad{-0.48}$, $n=99$).
\end{itemize}
\column{0.32\textwidth}
\centering\scriptsize \textbf{Within-city road-segment directionality}\par
\vspace{0.15em}
\includegraphics[width=\linewidth,height=0.40\textheight,keepaspectratio]{pat_bj_centerline_map.png}
\vspace{-0.15em}
\begin{itemize}\scriptsize
  \item Median \SAR\ is \good{0.884} among 3{,}670 Beijing two-way centerlines.
  \item \SAR\ $\leq1/3$ flags \good{179} candidate corridors, about \good{213 km}.
\end{itemize}
\end{columns}
\end{frame}

\begin{frame}[label=jobhousing]{Directionality Arises from Jobs--Housing Mismatch}
\framesubtitle{Square-grid commuting sample.}
\begin{columns}[T,onlytextwidth]
\column{0.43\textwidth}
\begin{itemize}\small\setlength\itemsep{0.75em}
  \item \textbf{Separation.} Jobs cluster more than residences: Duncan $D=\good{0.26}$, far above a random-allocation floor of $0.004$. \app{appduncan}
  \item \textbf{Direction.} AM flow is radial-inward: \good{74.5\%} of commuting trips head to the center, with inbound/reverse $\good{2.9\times}$. \app{appinbound}
\end{itemize}
\column{0.55\textwidth}
\centering\scriptsize Jobs--housing separation on the study grid\\
\includegraphics[width=\linewidth,height=0.64\textheight,keepaspectratio]{pat_jobs_housing_balance.png}
\end{columns}
\end{frame}

% =====================================================================
\section{Model}
% =====================================================================
% Original model section, preserved from tidalraw.tex.

\begin{frame}{Model: Directed Network and Costs}
\begin{itemize}\small
  \item Locations $i,j\in\mathcal N$ host a fixed mass $\bar L$ of residents and jobs.
  \item A physical road is represented by reciprocal directed links, $(k,l)$ and $(l,k)$.
  \begin{itemize}\footnotesize
    \item $m_{kl}^{ff}$ is free-flow physical time; $n_{kl}$ is directed lane capacity.
    \item Realized physical time $m_{kl}$ is endogenous and may differ from $m_{lk}$.
  \end{itemize}
  \item Keep three cost objects distinct:
  \vspace{-0.35em}
  \begin{equation}
    \underbrace{m_{kl}}_{\text{physical minutes}}
    \quad\Longrightarrow\quad
    \underbrace{t_{kl}=m_{kl}^{1/\theta}}_{\text{link iceberg cost}}
    \quad\Longrightarrow\quad
    \underbrace{\tau_{ij}}_{\text{bilateral commuting cost}}.
  \end{equation}
  \vspace{-0.6em}
  \item A tidal-lane policy reallocates $n_{kl}$ across reciprocal directions while holding total road width fixed.
\end{itemize}
\end{frame}
\note{\textbf{Talk track, 40 seconds.} Begin with the economic unit: locations are connected by directed links. Physical minutes, link iceberg costs, and bilateral commuting costs are separate objects. The policy changes directed capacity, while traffic determines realized physical time.}

\begin{frame}{Residence, Workplace, and Route Choice}
\begin{itemize}\small
  \item An individual $\omega$ living in $i$, working in $j$, and taking route $r\in\R_{ij}$ obtains
  \vspace{-0.5em}\begin{footnotesize}\begin{equation}
    V_{ij,r}(\omega) = \left( \frac{u_i w_j}{C_{ij,r}} \right)\varepsilon_{ij,r}(\omega),
    \qquad
    C_{ij,r} = \prod_{\ell=1}^{K_r} t_{r_{\ell-1},r_\ell}
  \vspace{-1.25em}\end{equation}\end{footnotesize}
  \begin{itemize}
    \item $u_i$ is amenity; $w_j$ is wage; $\epsilon$ is i.i.d. Fr\'echet.
    \item $C_{ij,r}$ is the route iceberg cost; route $r=[r_0,\dots,r_{K_r}]$, with $r_0=i$ and $r_{K_r}=j$.
  \end{itemize}
  \vspace{0.25em}
  \item Choice probabilities are
  \vspace{-0.5em}\begin{equation}\label{eq:pi}
    \pi_{ij,r} = \frac{\left(u_i w_j/C_{ij,r}\right)^\theta}{\sum_{a}\sum_b\sum_{s\in\R_{ab}}\left(u_a w_b/C_{ab,s}\right)^\theta}
  \end{equation}
\end{itemize}
\end{frame}
\note{\textbf{Talk track, 55 seconds.} The choice object is the triple $(i,j,r)$, with the origin--destination pair as one component. Fr\'echet heterogeneity converts utility into tractable choice shares. Relative to a bilateral-cost QSM, the route is inside the choice set, so changes on one directed segment propagate through route substitution.}

\begin{frame}[label=bilateral]{Bilateral Commuting Costs and Flows}
\begin{itemize}\footnotesize
  \item Route-inclusive bilateral cost comes from the Fr\'echet inclusive value:
  \vspace{-0.6em}\begin{equation}
    \tau_{ij}^{-\theta} = \sum_{r\in\R_{ij}} C_{ij,r}^{-\theta}
  \end{equation}
  \vspace{-0.8em}
  \item The same route aggregator has a network representation:
  \vspace{-0.55em}\begin{equation}\label{eq:inverse}
    A_{kl}=t_{kl}^{-\theta},\qquad B=(I-A)^{-1},\qquad \tau_{ij}=B_{ij}^{-1/\theta}.
  \end{equation}
  \vspace{-0.65em}
  \item With residence and workplace marginals, flows take a doubly constrained gravity form:
  \vspace{-0.55em}\begin{equation}\label{eq:gravity}
    L_{ij} = \sum_{r\in\R_{ij}}L_{ij,r} = \tau_{ij}^{-\theta} \left(\frac{L_i^R}{\Pi_i}\right) \left(\frac{L_j^F}{P_j}\right), \,\, L_i^R=\sum_jL_{ij}, \,\, L_j^F=\sum_iL_{ij}
  \vspace{-0.9em}\end{equation}
  \begin{itemize}
    \item $\Pi_i$ and $P_j$ balance the observed residence and workplace marginals.
    \item Directed link costs imply $\tau_{ij}\neq\tau_{ji}$ in general.
  \end{itemize}
\end{itemize}
\end{frame}
\note{\textbf{Talk track, 60 seconds.} Summing over routes converts route choice into one bilateral cost. The matrix inverse performs the same aggregation over feasible path lengths. Gravity determines OD commuter mass, and the next slide maps that mass back to directed segments. This is the bridge from OD gravity to segment-level congestion.}

\begin{frame}{Amenity, Productivity, and Location Choices}
\begin{itemize}\small
  \item Local outcomes combine fundamentals with endogenous density:
  \vspace{-0.5em}\begin{equation}\label{eq:agglom}
    w_j = A_j = \bar A_j\left(L_j^F\right)^\alpha, \qquad u_i=\bar u_i\left(L_i^R\right)^\beta
  \end{equation}
  \vspace{-0.4em}
  \item Residence ($R$) and workplace ($F$) shares satisfy
  \vspace{-0.2em}
  \begin{align}
  (l_i^R)^{1-\beta\theta} & = \chi \sum_j\tau_{ij}^{-\theta}\,\bar u_i^{\theta}\,\bar A_j^{\theta}\,(l_j^F)^{\alpha\theta} \label{eq:clear1} \\
  (l_j^F)^{1-\alpha\theta} & = \chi \sum_i\tau_{ij}^{-\theta}\,\bar u_i^{\theta}\,\bar A_j^{\theta}\,(l_i^R)^{\beta\theta} \label{eq:clear2}
  \end{align}
  \vspace{-0.7em}
  \item $\alpha>0$ governs workplace agglomeration; $\beta<0$ governs residential crowding; $\chi$ normalizes shares to one.
\end{itemize}
\end{frame}
\note{\textbf{Talk track, 55 seconds.} Fundamentals are $\bar u_i$ and $\bar A_j$. Density changes realized amenities and productivity through $\beta$ and $\alpha$. The two fixed-point equations determine long-run sorting once bilateral commuting costs change. Do not derive the equations; interpret the residence side and the workplace side once each.}

\begin{frame}[label=equilibrium]{Link Traffic, Congestion, and Equilibrium}
\begin{itemize}\footnotesize
  \item For a commuter traveling from $i$ to $j$, $\kappa^{kl}_{ij}$ is expected use of directed link $(k,l)$:
  \vspace{-0.85em}\begin{equation}
    \kappa^{kl}_{ij} = \left( \frac{\tau_{ij}} {\tau_{ik}\,t_{kl}\,\tau_{lj}} \right)^\theta
  \vspace{-0.65em}\end{equation}
  \item Model-implied directed link load aggregates OD mass across routes:
  \vspace{-0.55em}
  \begin{equation}\label{eq:Xi_c}
  \Xi_{kl} = \sum_i\sum_j \kappa^{kl}_{ij}L_{ij}
  \end{equation}
  \vspace{-0.75em}
  \item Traffic per lane determines physical time through an increasing speed--flow relation:
  \vspace{-0.55em}\begin{equation}\label{eq:c_Xi}
    m_{kl}=\mathcal M_{kl}\!\left(\frac{\Xi_{kl}}{n_{kl}}\right),
    \qquad \mathcal M_{kl}'>0.
  \end{equation}
  \vspace{-0.75em}
  \item Equilibrium jointly determines
  \vspace{-0.45em}\begin{equation}
    \{L^R,L^F,L_{ij},\tau_{ij},m_{kl},\Xi_{kl},\Pi_i,P_j,\bar W\}.
  \end{equation}
  \vspace{-0.75em}
  \item Location choice, route traffic, physical time, and bilateral costs are mutually consistent. \app{appexist}
\end{itemize}
\end{frame}
\note{\textbf{Talk track, 90 seconds.} Interpret $\kappa_{ij}^{kl}$ as expected passages over a directed segment for one $i$ to $j$ commuter. Multiplying by OD flows and summing gives $\Xi_{kl}$, a model-implied link load that does not rely on observed traffic counts. Congestion raises physical time, which changes route-inclusive bilateral costs and feeds back into location choice. This feedback is the endogenous-congestion component inherited from the Allen--Arkolakis framework; the empirical anchor is introduced after the data measurement slides.}

% =====================================================================
\section{Data and Estimation}\label{sec:quant}
% =====================================================================
\begin{frame}[label=data]{Model Data: From Observations to Spatial Objects}
\begin{columns}[T,onlytextwidth]
\column{0.54\textwidth}
\begin{itemize}\footnotesize\setlength\itemsep{0.45em}
  \item \textbf{Residence and workplace.} Baidu 2022 home and work points are assigned to $3\times3$ km cells and aggregated to $(L_i^R,L_j^F)$.
  \item \textbf{OD flows.} The full commuting matrix supports descriptive facts and validation; the baseline inversion targets the residence and workplace marginals.
  \item \textbf{Directed network.} Road geometry, functional class, lanes, and directional speed records are mapped to adjacent grid links.
  \item \textbf{Analysis sample.} The largest connected component is retained after a detour filter removes implausible links.
\end{itemize}
\column{0.44\textwidth}
\centering\scriptsize Directional speed records\par
\includegraphics[width=\linewidth,height=0.66\textheight,keepaspectratio]{speed_eg.png}
\end{columns}
\end{frame}
\note{\textbf{Talk track, 45 seconds.} Separate observed objects from model-implied objects. Residence and workplace marginals discipline the inversion. OD data support the descriptive section and validation, while route and segment loads are generated inside the model.}

\begin{frame}{Directed Link Times Are Constructed from Segment Speeds}
\begin{columns}[T,onlytextwidth]
\column{0.49\textwidth}
\begin{itemize}\footnotesize\setlength\itemsep{0.35em}
  \item \textbf{Observed AM time.} Weekday speeds from 07:00--09:00 deliver $v_a^{obs}$.
  \item \textbf{Free-flow time.} Night speeds from 22:00--06:00 deliver $v_a^{ff}$.
  \item \textbf{Pair-specific routing.} For each adjacent grid pair $(k,l)$, a road-network path $p_{kl}$ is selected and segment minutes are aggregated:
  \vspace{-0.35em}
  \begin{equation}
    m_{kl}^{q}=\sum_{a\in p_{kl}}\frac{\ell_a}{v_a^{q}},
    \qquad q\in\{obs,ff\}.
  \end{equation}
  \vspace{-0.65em}
  \item Direction is retained, so $m_{kl}^{q}$ need not equal $m_{lk}^{q}$.
\end{itemize}
\column{0.49\textwidth}
\centering\scriptsize AM physical time on adjacent grid links\par
\includegraphics[width=\linewidth,height=0.67\textheight,keepaspectratio]{app_edgeproj_adjacent_tau.png}
\end{columns}
\end{frame}
\note{\textbf{Talk track, 50 seconds.} Segment speeds are first converted into physical minutes. Pair-specific routing then maps the detailed road network to directed adjacent-grid links. The model route aggregator operates on these directed links; it does not replace this measurement step with a single shortest-path OD cost.}

\begin{frame}[label=estparams]{Parameters and Estimation Targets}
\small
\begin{columns}[T,onlytextwidth]
\column{0.58\textwidth}
\scriptsize
\begin{tabularx}{\linewidth}{@{}p{0.18\linewidth}p{0.16\linewidth}X@{}}
\toprule
Parameter & Value & Economic role \\
\midrule
$\theta$ & $6.83$ & response to commuting costs \\
$\alpha$ & $+0.05$ & workplace agglomeration \\
$\beta$ & $-0.10$ & residential crowding \\
$\lambda_{\mathrm{phys}}$ & $1.1$ & physical-time congestion elasticity \\
\bottomrule
\end{tabularx}
\column{0.40\textwidth}
\begin{itemize}\footnotesize
  \item \textbf{Matched baseline objects:} $L_i^R$ and $L_j^F$.
  \item \textbf{Measured network objects:} $m_{kl}^{obs}$, $m_{kl}^{ff}$, $n_{kl}$, and $\ell_{kl}$.
  \item \textbf{Model-implied objects:} $L_{ij}$, route shares, and $\Xi_{kl}$.
\end{itemize}
\end{columns}
\vspace{0.45em}
{\footnotesize $\theta$, $\alpha$, $\beta$, and $\lambda_{\mathrm{phys}}$ are externally calibrated in the current specification. \app{appparams}}
\end{frame}

\begin{frame}{Estimation: Invert, Anchor, and Solve}
\begin{itemize}\small
  \item \textbf{Invert fundamentals.} Recover $(\bar u_i,\bar A_j)$ so that observed residence and workplace marginals form an equilibrium at $m_{kl}^{obs}$.
  \item \textbf{Construct the baseline link-load anchor.} $\Xi_{kl}^{0}$ is the model-implied directed link load evaluated at baseline allocation and observed directed times.
  \item \textbf{Update physical time around that operating point:}
  \vspace{-0.45em}\begin{equation}\label{eq:anchor}
    m_{kl}=\max\!\left\{
    m_{kl}^{ff},\;
    m_{kl}^{obs}
    \left[
    \frac{\Xi_{kl}}{\Xi_{kl}^{0}}
    \frac{n_{kl}^{0}}{n_{kl}}
    \right]^{\lambda_{\mathrm{phys}}}
    \right\}.
  \end{equation}
  \vspace{-0.65em}
  \item \textbf{Solve in levels.} Each scenario recomputes $(m,t,\tau,\Xi,L^R,L^F,\Pi,P,\bar W)$ with fundamentals fixed.
\end{itemize}
\vspace{-0.25em}
\begin{center}\footnotesize
$\Xi_{kl}^{0}$ is a baseline model-implied load, not an observed vehicle count.
\end{center}
\end{frame}
\note{\textbf{Talk track, 55 seconds.} The inversion first recovers the fundamentals required to rationalize the observed marginals. The anchor fixes the baseline operating point at observed physical time. Only the ratio $\Xi_{kl}/\Xi_{kl}^{0}$ enters, so the arbitrary scale of model-implied link load cancels. Counterfactuals change traffic per lane and the model solves the full fixed point in levels.}

\begin{frame}[label=fundamentals]{Baseline Fundamentals (inverted)}
\begin{columns}[T,onlytextwidth]
\column{0.5\textwidth}
\centering \scriptsize Residence amenity $\bar u_i$\\
\includegraphics[width=\linewidth,height=0.52\textheight,keepaspectratio]{fund_ubar.png}
\column{0.5\textwidth}
\centering \scriptsize Workplace productivity $\bar A_j$\\
\includegraphics[width=\linewidth,height=0.52\textheight,keepaspectratio]{fund_abar.png}
\end{columns}
\vspace{0.1em}
\begin{itemize}\footnotesize
  \item $\bar u,\bar A$ are inverted to match observed residence and workplace shares at observed travel times.
  \item $\mathrm{corr}(\bar A_j,L_j^F)=\good{0.67}$ vs.\ $\mathrm{corr}(\bar u_i,L_i^R)=0.30$: jobs pinned by productivity, residence is access-driven.
  \item Treated corridors land on \bad{high-flow employment bottlenecks} (TidalCandidateScore vs.\ slow-direction $\Xi_s$: corr.\ $\approx+0.91$ on selected top-100 pairs). \app{appfund}
\end{itemize}
\end{frame}
\note{\textbf{Talk track, 35 seconds.} These maps are recovered objects, not directly observed amenities or productivity. The inversion matches baseline residence and workplace marginals at observed directed times. Use the maps to establish the spatial fundamentals that remain fixed in the tidal-lane counterfactuals.}

% =====================================================================
\section{Counterfactuals}
% =====================================================================
\begin{frame}{What Symmetry Misses}
\framesubtitle{The directional and symmetric measurement systems are re-estimated separately.}
\small
\begin{center}
\scriptsize
\begin{tabularx}{0.94\textwidth}{@{}p{0.18\textwidth}p{0.25\textwidth}p{0.24\textwidth}X@{}}
\toprule
World & Link-time measurement & Fundamental inversion & Equilibrium \\
\midrule
Directional & Directed $m_{kl}^{obs}$ and $m_{kl}^{ff}$ & Re-estimate $(\bar u,\bar A)$ & Re-solve anchored GE \\
Symmetric & Pairwise geometric mean & Re-estimate $(\bar u,\bar A)$ & Re-solve anchored GE \\
\bottomrule
\end{tabularx}
\end{center}
\vspace{0.25em}
\begin{itemize}\footnotesize
  \item \textbf{Primary estimand.} Changes in recovered fundamentals, market access, and equilibrium spatial gradients induced by imposing reciprocal symmetry.
  \item \textbf{Distributional diagnostics.} Compare dispersion, CBD gradients, and dissimilarity in $\bar u$, $\bar A$, bilateral costs, and model-implied link loads.
  \item \textbf{Policy relevance.} Re-evaluate the tidal-lane counterfactual under each re-estimated measurement system.
\end{itemize}
\vspace{0.3em}
\begin{center}
\fcolorbox{gray!50}{gray!8}{\parbox{0.77\textwidth}{\centering\footnotesize
[Full re-estimation results pending.] The fixed-fundamentals diagnostic is retained in the appendix.}}
\end{center}
\end{frame}

\begin{frame}{Where the Tidal Policy Lands}
\framesubtitle{Top-100 corridors by TidalCandidateScore (TCS).}
\begin{columns}[T,totalwidth=0.96\textwidth]
\column{0.42\textwidth}
\begin{itemize}\footnotesize
  \item For each reciprocal pair, label the slower AM direction $s$ and the faster direction $f$.
  \item Reallocate one lane while holding total road width fixed:
  \vspace{-0.35em}
  \[
    n_s'=n_s+1,\qquad n_f'=n_f-1.
  \]
  \vspace{-0.75em}
  \item \textbf{TCS} is the slow-minus-fast congestion-delay mass:
  \vspace{-0.35em}\[
    \begin{aligned}
    \mathrm{TCS}={}&\Xi_s(m_s^{obs}-m_s^{ff})\\
    &-\Xi_f(m_f^{obs}-m_f^{ff}).
    \end{aligned}
  \]
  \vspace{-0.65em}
  \item Red edges receive the lane among the top-100 reciprocal pairs.
\end{itemize}
\column{0.54\textwidth}
\centering
\vspace{-0.2em}
\includegraphics[width=\linewidth,height=0.78\textheight,keepaspectratio]{map_treated_overlay.png}
\end{columns}
\end{frame}
\note{\textbf{Talk track, 60 seconds.} The intervention is capacity reallocation, not capacity expansion. One lane moves from the faster AM direction to the slower direction on each selected reciprocal corridor. TCS ranks corridors by the difference in baseline flow times excess physical travel time. The map identifies the lane-receiving directions.}

\begin{frame}{Short Run vs.\ Long Run}
\small
\begin{itemize}
  \item \textbf{Short run (ST):} pin residence $L^R$ and employment $L^F$ at baseline; route traffic, physical link times $m_{kl}$, and bilateral costs $\tau_{ij}$ re-equilibrate.
  \item \textbf{Long run (LT):} allow $L^R$, $L^F$, and OD shares to sort, with congestion feedback.
  \item $s_{ij}^h$ is the OD commuter share and $\bar\tau^h=\sum_{ij}s_{ij}^h\tau_{ij}^h$ for $h\in\{0,ST,LT\}$.
\end{itemize}
\vspace{-0.55em}
{\footnotesize
  \begin{equation}
    \Delta\bar\tau^{LT}=
    \underbrace{\textstyle\sum_{ij}s_{ij}^0(\tau_{ij}^{ST}\!-\tau_{ij}^0)}_{\text{ST direct}}
    +\underbrace{\textstyle\sum_{ij}s_{ij}^0(\tau_{ij}^{LT}\!-\tau_{ij}^{ST})}_{\text{GE wedge}}
    +\underbrace{\textstyle\sum_{ij}(s_{ij}^{LT}\!-s_{ij}^0)\tau_{ij}^{LT}}_{\text{sorting}}.
  \end{equation}
}
\vspace{-0.8em}
\begin{itemize}\scriptsize
  \item The first two terms hold baseline OD weights fixed; the third term isolates reweighting across OD pairs.
\end{itemize}
\end{frame}
\note{\textbf{Talk track, 60 seconds.} The short run permits traffic and congestion to adjust but fixes residence and workplace shares. The long run also permits sorting. The decomposition separates the direct congestion response, the additional network GE response, and the change in OD weights.}

\begin{frame}[label=result2]{Results: Short-Run Gains and Long-Run Adjustment}
\begin{columns}[T,onlytextwidth]
\column{0.54\textwidth}
\centering\scriptsize
\setlength{\tabcolsep}{4pt}
\begin{tabular}{lcc}
\toprule
Channel & sec./commuter & person-min/day \\
\midrule
ST direct & \good{$-1.63$} & \good{$-2{,}400$} \\
$+$ network GE & $+0.99$ & $+1{,}456$ \\
$+$ sorting & \bad{$+0.71$} & \bad{$+1{,}039$} \\
\midrule
\textbf{Long-run total} & \textbf{$+0.06$} & \textbf{$+95$} \\
\bottomrule
\end{tabular}
\vspace{0.55em}
\begin{itemize}\footnotesize
  \item \good{56.4\%} of baseline OD commuter mass is exposed to treated corridors.
  \item Treated slow-direction links save \good{$20.3\%$} in ST and \good{$13.6\%$} in LT.
\end{itemize}
\column{0.39\textwidth}
\begin{center}
\small\textbf{Expected-utility welfare}\par
\vspace{0.5em}
\begin{tabular}{lc}
\toprule
Equilibrium & $\Delta\bar W$ \\
\midrule
Short run & \good{$+0.87\%$} \\
Long run & \good{$+0.32\%$} \\
\bottomrule
\end{tabular}
\end{center}
\vspace{0.65em}
\begin{itemize}\footnotesize
  \item Congestion feedback and sorting absorb most of the average route-time gain.
  \item Welfare remains positive because the choice set includes the full residence--workplace opportunity set.
  \item Mode switching is outside the current accounting.
\end{itemize}
\end{columns}
\end{frame}
\note{\textbf{Talk track, 90 seconds.} Read the left table as a sequence. The policy produces an immediate time saving, network congestion feedback offsets part of it, and long-run sorting offsets the remainder. The expected-utility welfare measure remains positive in both horizons because access to the full residence--workplace opportunity set changes even when the OD-weighted time total returns close to baseline.}

\begin{frame}{Long-Run Sorting}
\begin{columns}[T,onlytextwidth]
\column{0.35\textwidth}
\begin{itemize}\footnotesize
  \item Totals are fixed at $\bar L$; maps report demeaned share changes, not citywide growth.
  \item Positive cells gain share relative to baseline; black circles mark treated-corridor endpoint nodes.
  \item \textbf{Residential shares} rise near treated endpoints because lower OD costs raise residential market access.
  \item \textbf{Employment gains} are concentrated in a smaller set of red cells.
\end{itemize}
\column{0.63\textwidth}
\centering
{\scriptsize
\begin{tabular}{lcccc}
\toprule
 & \multicolumn{2}{c}{Treated endpoints} & \multicolumn{2}{c}{Rest} \\
 & $\ell^R$ & $\ell^F$ & $\ell^R$ & $\ell^F$ \\
\midrule
Mean $\Delta\log$ & \good{$+0.017$} & $-0.006$ & \bad{$-0.012$} & \bad{$-0.016$} \\
\bottomrule
\end{tabular}
}

\vspace{1em}
\begin{minipage}[t]{0.49\linewidth}
\centering \scriptsize demeaned $\Delta\log\ell^R$\\
\includegraphics[width=\linewidth,height=0.45\textheight,keepaspectratio]{sort_lR.png}
\end{minipage}\hfill
\begin{minipage}[t]{0.49\linewidth}
\centering \scriptsize demeaned $\Delta\log\ell^F$\\
\includegraphics[width=\linewidth,height=0.45\textheight,keepaspectratio]{sort_lF.png}
\end{minipage}
\end{columns}
\end{frame}

\begin{frame}[label=result4]{Transportation Policy Comparison}
\framesubtitle{Endogenous GE under anchored congestion and long-run sorting.}
\begin{center}\scriptsize
\begin{tabular}{lccccc}
\toprule
Policy & ST $\bar W$ & LT $\bar W$ & LT route cost & net lane-km & cost/lane-km \\
 &  &  & (person-min/day) &  &  \\
\midrule
Tidal reallocation & $+0.87\%$ & $+0.32\%$ & $+95$ & \good{\textbf{0}} & near zero capex \\
Road Expansion & $+1.26\%$ & $+0.44\%$ & $+75$ & 1{,}140 & $\approx$ CNY 5m \\
\bottomrule
\end{tabular}
\end{center}
\vspace{0.3em}
\begin{itemize}\footnotesize
  \item Tidal reallocation lifts LT welfare $+0.32\%$ at zero net lane-km.
  \item Delivering comparable long-run welfare by building capacity takes about $1{,}140$ lane-km: Road Expansion adds lanes where baseline flow$\times$delay is highest and reaches $+0.44\%$.
  \item Urban arterial widening runs on the order of CNY 5m per lane-km (roughly CNY 30m/km for a six-lane arterial), so the expansion benchmark is order CNY 6bn; tidal reallocation avoids this outlay. \app{apppolicy}
\end{itemize}
\end{frame}

\begin{frame}{Conclusion}
\begin{enumerate}
  \item \textbf{Directional congestion is a widespread urban feature.}
  \begin{itemize}\footnotesize
    \item Across the top 100 congested cities, flow-weighted median \SAR\ is \good{0.878}; the same asymmetry persists on Beijing road segments.
  \end{itemize}
  \item \textbf{Directional measurement belongs inside the structural inversion.}
  \begin{itemize}\footnotesize
    \item The full comparison re-estimates fundamentals and equilibrium gradients under the directional and reciprocal-symmetric cost surfaces.
  \end{itemize}
  \item \textbf{Tidal lanes deliver welfare gains through directed capacity reallocation.}
  \begin{itemize}\footnotesize
    \item ST route cost falls by \good{$1.63$ s/commuter} and $\bar W$ rises \good{$+0.87\%$}; LT sorting absorbs part of the time gain, but $\bar W$ remains \good{$+0.32\%$}.
    \item Supply-side policies with comparable LT welfare require about \bad{$1{,}100$ lane-km}; tidal reallocation uses \good{zero net lane-km}.
  \end{itemize}
\end{enumerate}
\end{frame}

% =====================================================================
\appendix
\section{Appendix}
% =====================================================================
\begin{frame}[label=jobhousingnational]{Appendix: Jobs--Housing Mismatch Across Major Cities}
\framesubtitle{Township commuting flows, top 20 cities by inter-town commuter mass.}
\begin{itemize}\footnotesize\setlength\itemsep{0.25em}
  \item Core towns cover the first 30\% of workplace-side commuter mass in each city.
  \item In all top-20 cities, $P\to C$ exceeds $C\to P$; the flow-weighted gap is \good{4.5 pp}.
  \item The same gap is \good{0.9 pp} across all cities, so the major-city sample carries much of the directionality pattern.
\end{itemize}
\vspace{-0.15em}
\centering
\makebox[\linewidth][c]{\includegraphics[width=1.08\linewidth,height=0.66\textheight,keepaspectratio]{pat_jobs_housing_top20_in_out}}
\end{frame}

\begin{frame}[label=jobhousingasym]{Appendix: Jobs--Housing Mismatch and Speed Asymmetry}
\framesubtitle{City-level comparison with national township-pair AM \SAR.}
\begin{itemize}\scriptsize\setlength\itemsep{0.15em}
  \item Lower AM \SAR\ means stronger pair-level speed asymmetry; the top-20 cross-city relationship is weak and retained as a diagnostic.
  \item The figures use the same employment-core definition and compare city-level $P\to C$ imbalance with AM speed asymmetry.
\end{itemize}
\vspace{0.2em}
\centering
\makebox[\linewidth][c]{\includegraphics[width=1.02\linewidth,height=0.62\textheight,keepaspectratio]{pat_jobs_housing_asym_scatter}}
\end{frame}

\begin{frame}[label=jobhousingdumbbell]{Appendix: Directional Commuting Shares by City SAR}
\framesubtitle{Top 20 cities, same employment-core definition.}
\begin{itemize}\scriptsize\setlength\itemsep{0.15em}
  \item Each city contributes two points at its AM \SAR: the $P\to C$ share and the $C\to P$ share.
  \item All top-20 city gaps are positive, reflecting higher center-bound commuting shares in major cities.
\end{itemize}
\vspace{-0.35em}
\centering
\includegraphics[width=1.05\linewidth,height=0.72\textheight,keepaspectratio]{pat_jobs_housing_asym_dumbbell}
\end{frame}

\begin{frame}[label=appsymfixed]{Appendix: Fixed-Fundamentals Symmetry Diagnostic}
\framesubtitle{This estimand differs from the full re-estimation in the main deck.}
\footnotesize
Hold the network, lanes, inverted fundamentals, parameters, solver, and welfare formula fixed; replace reciprocal-pair $m_{kl}^{obs}$ and $m_{kl}^{ff}$ by pairwise geometric means.
\vspace{0.45em}
\begin{center}
\scriptsize
\begin{tabular}{lcccc}
\toprule
 & $\Delta\bar W$ & Residence & Workplace & OD cost \\
 &  & dissim. & dissim. & $\Delta$ \\
\midrule
Symmetric vs.\ directed & \good{$+1.10\%$} & $3.36\%$ & $6.15\%$ & $+0.10\%$ \\
\bottomrule
\end{tabular}
\end{center}
\vspace{0.35em}
\begin{itemize}\footnotesize
  \item The experiment changes the measured cost surface while holding fundamentals fixed.
  \item It is not a comparison of two fully re-estimated models.
\end{itemize}
\end{frame}

\begin{frame}[label=approute]{Appendix: Route-Specific Traversal Cost vs.\ Residence-Workplace Commuting Cost}
\begin{itemize}\small
  \item Define the directed adjacency matrix $A=[a_{kl}]$ by
  \vspace{-0.5em}\begin{equation}
    a_{kl} = t_{kl}^{-\theta}
  \end{equation}
  \item Summing over all feasible route lengths gives
  \vspace{-0.5em}\begin{equation}
    B = \sum_{K=0}^{\infty}A^K = (I-A)^{-1}
  \end{equation}
  \item The bilateral commuting-cost aggregator can be written as
  \vspace{-0.5em}\begin{equation}\label{eq:inverse_app}
    \tau_{ij} = B_{ij}^{-1/\theta}
  \end{equation}
  \vspace{-0.5em}
  \item The implementation applies a common link-cost rescaling $\zeta$ to keep the spectral radius of $A$ below one. This scalar is distinct from $\kappa_{ij}^{kl}$, the expected use of link $(k,l)$ by an $i\to j$ commuter.
\end{itemize}
\navback{bilateral}
\end{frame}

\begin{frame}[label=appfund]{Appendix: Estimation Algorithm}
\footnotesize
\textbf{Baseline inversion}
\begin{itemize}
  \item Treat observed residence and workplace marginals as the baseline equilibrium.
  \item Recover $(\bar u_i^\theta,\bar A_j^\theta)$ by alternating updates with geometric-mean normalization.
  \item Match $(L_i^R,L_j^F)$, not the full bilateral OD matrix; OD flows are model-implied by gravity.
\end{itemize}
\vspace{0.25em}
\textbf{Counterfactual solve}
\begin{itemize}
  \item Given fundamentals and parameters, solve $(L^R,L^F)$, market access, model-implied link load $\Xi$, and physical times jointly.
  \item Anchored congestion updates physical time around the baseline model-implied load and observed lanes:
  \vspace{-0.3em}
  \[
    m_{kl}=\max\{m_{kl}^{ff},\,m_{kl}^{obs}[(\Xi_{kl}/\Xi_{kl}^{0})(n_{kl}^{0}/n_{kl})]^{\lambda_{\mathrm{phys}}}\}.
  \]
  \vspace{-0.8em}
  \item Outer loop updates $m_{kl}$; inner loop alternates location shares and market access until convergence.
\end{itemize}
\navback{fundamentals}
\end{frame}

\begin{frame}[label=appduncan]{Appendix: Jobs--Housing Dissimilarity Benchmarks}
\framesubtitle{Duncan $D=\tfrac12\sum_g\lvert j_g/J-r_g/R\rvert$ ($j,r=$ jobs, residents per unit $g$); aspatial, $0$ balanced $\to1$ separated.}
\begin{columns}[T,onlytextwidth]
\column{0.60\textwidth}
\centering
\includegraphics[width=\linewidth,height=0.58\textheight,keepaspectratio]{app_duncan_cross_city.png}
\column{0.38\textwidth}
\begin{itemize}\footnotesize
  \item \textbf{Cross-city.} Township $D$, top-100 cities (residents $=$ origin, jobs $=$ dest mass, 2021 flows): Beijing $\good{0.20}$ ranks \good{1/100} (median $0.09$).
  \item \textbf{Caveat.} $D$ is aspatial, rising with size (corr $0.50$; $0.21$ with unit count): Beijing's lead is not a unit-count artefact, but $D$ ignores distance.
\end{itemize}
\end{columns}
\navback{jobhousing}
\end{frame}

\begin{frame}[label=appinbound]{Appendix: Radial AM Inbound--Reverse Asymmetry}
\framesubtitle{Road-mode OD; \emph{inbound} $=$ home $\geq15$ km from the job-weighted center \& work $\leq10$ km; \emph{reverse} $=$ mirror.}
\begin{columns}[T,onlytextwidth]
\column{0.62\textwidth}
\centering
\includegraphics[width=\linewidth,height=0.66\textheight,keepaspectratio]{cityShenzhen2.jpg}
\column{0.36\textwidth}
\begin{itemize}\footnotesize
  \item Each core-directed trip is binned into one of 16 $22.5^\circ$ sectors by \emph{home} bearing.
  \item Pooled inbound share \good{$74.5\%$}; inbound/reverse $\good{2.9\times}$; \good{11/16} sectors $>0.70$ (min $0.57$, median $0.77$, max $0.86$).
  \item Non-circular: inbound and reverse share one center, so the ratio is a genuine imbalance, not an artefact of the job-cluster construction.
\end{itemize}
\end{columns}
\navback{jobhousing}
\end{frame}

\begin{frame}[label=appmodelasym]{Appendix: Model-Implied Directional Asymmetry}
\framesubtitle{Baseline directed traffic $\Xi$: ratio of the two directions per link (model output, not raw data).}
\setlength{\columnsep}{0.005\textwidth}
\begin{columns}[T,onlytextwidth]
\column{0.53\textwidth}
\centering
\includegraphics[width=\linewidth,height=0.70\textheight,keepaspectratio]{map_dir_asym.png}
\column{0.45\textwidth}
\begin{itemize}\footnotesize
  \item Color is $\log_{10}(\Xi_{\mathrm{high}}/\Xi_{\mathrm{low}})$; width is total reciprocal flow.
  \item Median two-way imbalance is \bad{$2.9\times$}; p90 is \bad{$20\times$}.
  \item High-imbalance links concentrate on the central high-flow spine.
  \item Object: equilibrium link load $\Xi$, not observed speed.
\end{itemize}
\end{columns}
\navback{equilibrium}
\end{frame}

\begin{frame}[label=appparams]{Appendix: Parameter Interpretation}
\begin{itemize}\small
  \item $\theta$ governs sensitivity to bilateral commuting costs. Higher $\theta$ makes households and jobs respond more sharply to directed cost changes.
  \item $\alpha>0$ is workplace agglomeration in $w_j=\bar A_j(L_j^F)^\alpha$. It reinforces employment concentration around locations that gain access.
  \item $\beta<0$ is residential crowding in $u_i=\bar u_i(L_i^R)^\beta$. It penalizes residential pile-up and dampens household concentration.
  \item $\lambda_{\mathrm{phys}}$ maps traffic per lane into directed physical time. Higher $\lambda_{\mathrm{phys}}$ raises the short-run congestion response to capacity reallocation.
  \item The corresponding iceberg-cost elasticity in the Allen--Arkolakis notation is $\lambda_{\mathrm{AA}}=\lambda_{\mathrm{phys}}/\theta$.
  \item Sorting reflects the interaction of all four forces: accessibility changes shift OD costs, $\alpha$ pulls jobs toward stronger access, and $\beta$ limits residential concentration.
\end{itemize}
\navback{estparams}
\end{frame}

\begin{frame}[label=appsens]{Appendix: Sensitivity Summary}
\framesubtitle{Final mainline rerun, TCS top-100, primitive $\bar W$, $\lambda_{\mathrm{phys}}=1.1$.}
\begin{columns}[T,onlytextwidth]
\column{0.5\textwidth}
\centering Vary $\alpha$ ($\beta=-0.10$)\\[2pt]
\scriptsize
\begin{tabular}{cccc}
\toprule
$\alpha$ & sort (s) & person-min/day & $\bar W$ LT \\
\midrule
$-0.12$ & $-3.63$ & $-6417$ & $+0.37\%$ \\
$0$     & $-1.33$ & $-2960$ & $+0.34\%$ \\
$+0.03$ & $-0.25$ & $-1343$ & $+0.33\%$ \\
$\mathbf{+0.05}$ & $\mathbf{+0.71}$ & $\mathbf{+95}$ & $\mathbf{+0.32\%}$ \\
$+0.10$ & $+4.88$ & $+6340$ & $+0.29\%$ \\
\bottomrule
\end{tabular}
\column{0.5\textwidth}
\centering Vary $\beta$ ($\alpha=+0.05$)\\[2pt]
\scriptsize
\begin{tabular}{cccc}
\toprule
$\beta$ & sort (s) & person-min/day & $\bar W$ LT \\
\midrule
$\mathbf{-0.10}$ & $\mathbf{+0.71}$ & $\mathbf{+95}$ & $\mathbf{+0.32\%}$ \\
$0$     & $-1.57$ & $-3173$ & $+0.31\%$ \\
$+0.05$ & $-4.64$ & $-7579$ & $+0.28\%$ \\
$+0.10$ & \multicolumn{3}{c}{\emph{not converged}} \\
\bottomrule
\end{tabular}
\end{columns}
\vspace{0.25em}
\begin{itemize}\footnotesize
  \item $\bar W$ LT remains positive in every converged $\alpha$--$\beta$ cell; route-time statistics move more because sorting changes OD weights.
  \item Detailed cells, $\theta$ grid, and congestion-band checks\app{appsensdetail}\app{appsens_theta}.
\end{itemize}
\navback{result2}
\end{frame}

\begin{frame}[label=appsensdetail]{Appendix: Sensitivity Details}
\scriptsize
\begin{columns}[T,onlytextwidth]
\column{0.40\textwidth}
\begin{flushright}
\setlength{\tabcolsep}{1.0pt}
\begin{tabular}{ccccc}
\toprule
$\alpha$ & $\beta$ & sort (s) & person-min/day & $\bar W$ LT \\
\midrule
$-0.12$ & $-0.10$ & $-3.63$ & $-6{,}417$ & $+0.37\%$ \\
$0.00$ & $-0.10$ & $-1.33$ & $-2{,}960$ & $+0.34\%$ \\
$0.03$ & $-0.10$ & $-0.25$ & $-1{,}343$ & $+0.33\%$ \\
$\mathbf{0.05}$ & $\mathbf{-0.10}$ & $\mathbf{+0.71}$ & $\mathbf{+95}$ & $\mathbf{+0.32\%}$ \\
$0.10$ & $-0.10$ & $+4.88$ & $+6{,}340$ & $+0.29\%$ \\
$0.05$ & $0.00$ & $-1.57$ & $-3{,}173$ & $+0.31\%$ \\
$0.05$ & $0.05$ & $-4.64$ & $-7{,}579$ & $+0.28\%$ \\
$0.05$ & $0.10$ & \multicolumn{3}{c}{\emph{not converged}} \\
\bottomrule
\end{tabular}
\end{flushright}
\column{0.58\textwidth}
\begin{itemize}\footnotesize
  \item ST direct effect is stable: about $-1.63$ s/commuter and $\bar W$ ST about $+0.87\%$.
  \item $\theta$ grid (\app{appsens_theta}): each value re-inverts fundamentals and GE; LT $\bar W$ gains decline with $\theta$ while signs stay positive.
  \item Congestion band: at $\lambda_{\mathrm{phys}}=0.57$, top-100 LT person-min/day is about $-1{,}168$ and $\bar W$ LT is $-0.078\%$; this is retained as a legacy lower-slope robustness, not the mainline.
  \item Selection robustness: the asymmetry-ratio rule is weaker than TCS selection and can flip LT $\bar W$ negative.
\end{itemize}
\end{columns}
\navback{appsens}
\end{frame}

\begin{frame}[label=appsens_theta]{Appendix: Theta Sensitivity}
\framesubtitle{TCS top-100; each $\theta$ re-inverts fundamentals and re-solves baseline + tidal ST/LT.}
\scriptsize
\begin{center}
\renewcommand{\arraystretch}{1.3}
\setlength{\tabcolsep}{12pt}
\begin{tabular}{ccccc}
\toprule
$\theta$ & ST $\bar W$ & LT $\bar W$ & res.\ dissim & job dissim \\
\midrule
$3$ & $+1.99\%$ & $+0.75\%$ & $1.11$ & $1.55$ \\
$5$ & $+1.19\%$ & $+0.45\%$ & $0.97$ & $1.67$ \\
$\mathbf{6.83}$ & $\mathbf{+0.87\%}$ & $\mathbf{+0.32\%}$ & $\mathbf{0.86}$ & $\mathbf{1.79}$ \\
$8$ & $+0.74\%$ & $+0.27\%$ & $0.80$ & $1.89$ \\
$10$ & $+0.59\%$ & $+0.21\%$ & $0.70$ & $2.07$ \\
\bottomrule
\end{tabular}
\end{center}
\vspace{0.4em}
\begin{itemize}\footnotesize
  \item LT $\bar W$ stays positive across the grid; magnitudes fall as $\theta$ rises (stronger cost sensitivity).
  \item Resident dissim falls and job dissim rises with $\theta$; baseline $\bar W$ levels are not comparable across $\theta$.
\end{itemize}
\note{Mainline $\theta=6.83$ matches the tidal diagnostics row; baseline $\bar W$ levels are not comparable across $\theta$.}
\navback{appsens}
\end{frame}

\begin{frame}[label=apppolicy]{Appendix: Policy Definitions and Selected Effects}
\scriptsize
\begin{center}
\resizebox{0.98\textwidth}{!}{%
\begin{tabular}{p{0.19\textwidth}p{0.39\textwidth}cccc}
\toprule
Policy & Capacity rule & ST $\bar W$ & LT $\bar W$ & lane-km & cost/lane-km \\
\midrule
Tidal reallocation & TCS top-100; add one lane in the slow direction and remove one from the reciprocal fast direction. & $+0.87\%$ & $+0.32\%$ & 0 & near zero capex \\
Road Expansion & Top 100 by baseline flow $\times$ excess travel time; add one lane to each direction. & $+1.26\%$ & $+0.44\%$ & 1{,}140 & $\approx$ CNY 5m \\
Uniform expansion & Scale all links by $+2.4\%$ so added lane-km matches the targeted-expansion budget. & $+0.21\%$ & $+0.10\%$ & 1{,}092 & $\approx$ CNY 5m \\
\bottomrule
\end{tabular}
}
\end{center}
\vspace{0.2em}
\begin{itemize}\footnotesize
  \item ST pins residence and workplace shares at baseline; LT re-solves residence--workplace sorting.
  \item Road Expansion asks where an engineer would add capacity using flow and delay only; uniform spreads the same lane-km budget across the whole network (targeting versus spreading).
  \item Cost benchmark uses 2023--24 municipal urban-arterial construction estimates ($\approx$ CNY 30m per km, roughly CNY 5m per widened lane-km, a single widened lane being a fraction of the full cross-section), not inter-city expressway costs; the targeted expansion is thus order CNY 6bn.
\end{itemize}
\navback{result4}
\end{frame}

\begin{frame}[label=resultelast]{Appendix: Targeting Tidal Lanes}
\framesubtitle{Per-link welfare elasticity from $d\bar W^\theta=dA_{ab}\phi_a\psi_b$.}
\begin{columns}[T,onlytextwidth]
\column{0.44\textwidth}
\centering
\includegraphics[width=\linewidth,height=0.66\textheight,keepaspectratio]{map_link_welfare_elast.png}
\column{0.53\textwidth}
\scriptsize
\textbf{Selection rule} (top-100):
\begin{tabular}{lcc}
\toprule
rule & LT $\bar W$ & person-min/day \\
\midrule
\good{welfare elast.} & \good{$+0.35\%$} & $-50$ \\
TCS score & $+0.32\%$ & $+95$ \\
speed only & $+0.13\%$ & $+1430$ \\
\bottomrule
\end{tabular}
\vspace{0.5em}
\begin{itemize}\footnotesize
  \item The local score decomposes into edge gain $\times$ upstream residence access $\times$ downstream workplace opportunity. \app{appelast}
  \item Flow matters: the speed-only rank correlation with welfare elasticity is \bad{$0.08$}; the flow$\times$delay correlation is $0.65$.
  \item Welfare-elasticity and TCS targeting select the high-flow spine; speed-only also selects slow, low-flow links.
\end{itemize}
\end{columns}
\end{frame}

\begin{frame}[label=appelast]{Appendix: Welfare-Elasticity Score}
\scriptsize
\begin{itemize}
  \item Define the directed-link accessibility matrix by $A_{ab}=t_{ab}^{-\theta}$ and the network multiplier by $B=(I-A)^{-1}$.
  \item A local gain on one directed edge uses the matrix-inverse differential:
  \[
    dB=B(dA)B.
  \]
  \item If only edge $a\to b$ changes, $dA$ has one nonzero entry, so the welfare-index numerator has the local derivative
  \[
    d\bar W^\theta=dA_{ab}\phi_a\psi_b.
  \]
  \item $dA_{ab}$ is the direct gain from one more lane; $\phi_a$ is upstream residence-access weight reaching node $a$; $\psi_b$ is downstream workplace-opportunity weight from node $b$.
  \item The score is a local ranking device. The reported LT $\bar W$ values still come from re-solving the endogenous GE counterfactual.
\end{itemize}
\navback{resultelast}
\end{frame}

\begin{frame}[label=nextsteps]{Appendix: Next Steps}
\small
\begin{itemize}
  \item \textbf{Reduced-form evidence on tidal lanes.}
  \begin{itemize}\footnotesize
    \item Record traffic at operational tidal-lane segments and recover directed vehicle counts and speeds.
    \item Estimate the local speed--flow relation around the lane-direction switch. \app{appfieldvisit}
  \end{itemize}
  \item \textbf{Structural parameters.}
  \begin{itemize}\footnotesize
    \item Tighten estimates of $\alpha$, $\beta$, $\theta$, and $\lambda_{\mathrm{phys}}$ as the data permit.
  \end{itemize}
  \item \textbf{External validity.}
  \begin{itemize}\footnotesize
    \item Extend the measurement and counterfactual pipeline to Beijing and Shenzhen, AM and PM peaks, and alternative grid resolutions.
  \end{itemize}
\end{itemize}
\end{frame}

\begin{frame}[label=appfieldvisit]{Appendix: Field Visit}
\framesubtitle{On-site survey of Shenzhen tidal-lane corridors.}
\vspace{-0.15em}
\setlength{\columnsep}{0.014\textwidth}
\begin{columns}[T,onlytextwidth]
\column{0.317\textwidth}
\centering
\includegraphics[width=\linewidth,height=0.305\textheight,keepaspectratio]{fieldvisit_crop_1.jpg}\par\vspace{0.24em}
\includegraphics[width=\linewidth,height=0.305\textheight,keepaspectratio]{fieldvisit_crop_2.jpg}
\column{0.317\textwidth}
\centering
\includegraphics[width=\linewidth,height=0.305\textheight,keepaspectratio]{fieldvisit_crop_5.jpg}\par\vspace{0.24em}
\includegraphics[width=\linewidth,height=0.305\textheight,keepaspectratio]{fieldvisit_crop_4.jpg}
\column{0.317\textwidth}
\centering
\includegraphics[width=\linewidth,height=0.305\textheight,keepaspectratio]{fieldvisit_crop_6.jpg}\par\vspace{0.24em}
\includegraphics[width=\linewidth,height=0.305\textheight,keepaspectratio]{fieldvisit_crop_7.jpg}
\end{columns}
\navback{nextsteps}
\end{frame}

\begin{frame}[label=appexist]{Appendix: Existence and Uniqueness}
\framesubtitle{Core Allen--Arkolakis system; sufficient conditions from \citet[Proposition~1]{allenarkolakis2022}.}
\small
\begin{itemize}
  \item \textbf{Existence.} On a strongly connected network a strictly positive equilibrium exists for any $\alpha,\beta$, $\theta>0$, and $\lambda_{\mathrm{AA}}\ge0$ (Brouwer fixed point on the transformed operator).
  \item \textbf{Uniqueness (sufficient), for $\alpha,\beta\in[-1,1]$:}
  \vspace{-0.15em}\[
    \alpha\le\tfrac12\!\left(\tfrac1\theta-\lambda_{\mathrm{AA}}\right)
    \qquad\text{and}\qquad
    \beta\le\tfrac12\!\left(\tfrac1\theta-\lambda_{\mathrm{AA}}\right).
  \]
  \item \textbf{Principle.} The proof bounds the equilibrium operator so that its elasticity map has spectral radius at most one \citep{allenarkolakisli2020,karlin1967}.
  \item \textbf{Mapping.} The empirical implementation uses $\lambda_{\mathrm{phys}}$ on physical time, with $\lambda_{\mathrm{AA}}=\lambda_{\mathrm{phys}}/\theta$ on link iceberg cost.
  \item \textbf{Scope.} The anchored implementation is solved numerically; the proposition applies to the core theoretical system.
  \item The condition is sufficient, not necessary. At $(\alpha,\beta)=(0.05,\,0.10)$ both spillovers turn agglomerative, so the iteration fails to converge.
\end{itemize}
\navback{equilibrium}
\end{frame}

\nobibliography{references}

\end{document}
